\section{Implementation}
\label{sec:implementation}

\tool is written in Python and has no cloud credentials in its analysis path.
It reads infrastructure as code and produces findings, so a deployment can be
analysed before it exists.

\subsection{From source to a state graph}

\texttt{extract\_graph} reads Terraform and produces the state shape the
analysis works on: principals, policies, roles, runtimes, and the arguments
that connect them. Two parsing decisions carry weight. When a declared name is an expression the
parser cannot resolve, it falls back to the resource label and records the
expression, which is what Section~\ref{sec:limitations} reports as
undecidable. And a runtime block carrying \texttt{count} or
\texttt{for\_each} creates one runtime per element, so counting blocks and
counting runtimes differ; Section~\ref{sec:results} reports both.

\subsection{One configuration, end to end}

Figure~\ref{fig:worked} carries a configuration from the corpus through the
whole analysis. The outcome is a blocked edge, not a chain.

\begin{figure*}[t]
\centering
\begin{minipage}[t]{0.325\textwidth}
\textbf{\small 1. What the deployment declares}
\begin{lstlisting}
resource "aws_bedrockagentcore_agent_runtime"
        "agent" {
  agent_runtime_name = "cloud_engineer_agent"
  role_arn           = aws_iam_role.agent_role.arn
  agent_runtime_artifact {
    container_configuration {
      container_uri = "${local.ecr}:latest" } }
}

resource "aws_iam_role" "agent_role" {
  name = "agent-runtime-role"
  assume_role_policy = jsonencode({
    Statement = [{
      Action    = "sts:AssumeRole"
      Effect    = "Allow"
      Principal = { Service =
        "bedrock-agentcore.amazonaws.com" } }] })
}

# policy attached to the CI role
Statement = [{
  Sid    = "IAMManagement"
  Effect = "Allow"
  Action = [ "iam:CreateRole", "iam:PassRole",
             "iam:PutRolePolicy", ... ]
  Resource = "*"
}]
\end{lstlisting}
\end{minipage}\hfill
\begin{minipage}[t]{0.325\textwidth}
\textbf{\small 2. The state the analysis reads}
\begin{lstlisting}
principals
  ci-deploy-role
  agent-runtime-role

agents
  agent   rho = agent-runtime-role

grants
  (ci-deploy-role, iam:PassRole,      *)
  (ci-deploy-role, iam:PutRolePolicy, *)
  (ci-deploy-role, iam:CreateRole,    *)

trust
  (bedrock-agentcore.amazonaws.com,
     agent-runtime-role)
  no condition key on the statement
\end{lstlisting}
\vspace{2pt}
{\scriptsize The wildcard is kept as a wildcard. It is matched against
\texttt{agent-runtime-role} when the edge rule is evaluated, not expanded into
a list of resources beforehand.}
\end{minipage}\hfill
\begin{minipage}[t]{0.325\textwidth}
\textbf{\small 3. What is reported}
\begin{lstlisting}
edge  ci-deploy-role --pass--> agent-runtime-role
  (*Eq.~\ref{eq:replace}*) conjunct 2  iam:PassRole over rho(agent)
                   HOLDS, Resource = "*"
                   unconditional
  (*Eq.~\ref{eq:replace}*) conjunct 1  UpdateAgentRuntime on agent
                   MISSING

verdict   the pass gate is open
chain     none
why       the configuration declares 1 agent,
          so there is no sibling to take over

reported as a precondition
not reported as a cross-agent escalation
\end{lstlisting}
\vspace{2pt}
{\scriptsize This is the \nPassRoleWide configuration of
Section~\ref{sec:results} where the condition that blocks the other
\nBlocked candidates is already absent.}
\end{minipage}
\caption{One configuration from the corpus, carried through the analysis.
Names are shortened and the repository is withheld until disclosure; the
structure, the wildcard, and the absent trust condition are as they appear in
the source. The finding is a precondition rather than a chain, and the reason
is recorded with it.}
\label{fig:worked}
\end{figure*}


Panel 1 is the source. A runtime names its execution role through a resource
reference rather than a literal ARN, the role's trust policy admits the
service with no condition key at all, and a policy attached to a continuous
integration role grants \iam{iam:PassRole} with \texttt{Resource = "*"}.

Panel 2 is the state. Three things in it decide the result. The runtime's
\texttt{role\_arn} is resolved through the reference to the role resource, so
$\rho(\mathit{agent})$ is known. The wildcard in the grant is carried as a
wildcard rather than expanded, so the match happens when the edge rule is
evaluated. The trust statement is recorded with the fact that it has no
condition, which is stronger than the vendor template of
Section~\ref{sec:background} and is what makes the trust side of the edge
unconditional here.

Panel 3 is the finding. Evaluating Equation~\ref{eq:replace} against this
state, the \iam{iam:PassRole} conjunct holds: the wildcard covers
\texttt{agent-runtime-role}, and no condition narrows it. The
\api{UpdateAgentRuntime} conjunct does not hold, because the role carries no
action on the agent platform. So the edge is blocked, the single missing
action is named, and no chain is reported.

The analysis then declines to escalate this further for a second reason, and
records it: the configuration declares one agent. Equation~\ref{eq:replace}
needs a runtime whose role differs from the caller's reach, and a deployment
with a single agent has no second role to reach. The output says both things,
so a reader can see that the pass gate being open and a cross-agent
escalation existing are different claims.

This is the \nPassRoleWide configuration that Section~\ref{sec:results}
reports separately from the \nBlocked blocked candidates. In the others the
missing conjunct is \iam{iam:PassRole}. Here it is the update action, and the
pass is already granted.

\subsection{Reaching}

\texttt{reach.py} computes $\auth$ and $\eff$ and applies
Equation~\ref{eq:esc}. CloudFormation and CDK are translated into the same state shape and analysed
by the same code, so a difference between corpora is a difference in the
deployments rather than in the analysis.

\subsection{CloudFormation}

\texttt{cfn\_to\_state} resolves the intrinsic functions that appear in role
and runtime arguments. An intrinsic it cannot resolve is recorded, and the
template is reported as undecidable rather than analysed. The function that
mattered most in practice was \api{Fn::Join}, because CDK writes nearly every
ARN as a join over a partition reference, and without it the CDK corpus could
not be analysed at all.

\subsection{CDK}

A CDK application builds its roles in code, so recovering the policies by
reading the source is a heuristic. We avoid the heuristic by synthesising:
each application is installed in a scratch directory, its entry point is run
with \texttt{CDK\_OUTDIR} set, and the CloudFormation it emits goes through
the same pipeline as the CloudFormation corpus. The CDK CLI is not needed,
because an application writes its templates when \texttt{App.synth()} is
called.

Failures to synthesise are reported by reason in Section~\ref{sec:results}.

\subsection{Scale}

Over the \nProjects analysed projects, which hold \nTerraformFiles Terraform
files, the whole corpus takes \corpusSeconds seconds: a median of
\medianSeconds seconds per template, \pNinetySeconds at the 90th percentile,
\pNinetyNineSeconds at the 99th, \maxSeconds at worst, and \peakMemoryMB\,MB
of peak child memory. That is cheap enough to run on every change rather than
periodically, which matters for a finding whose value is in review.
